\section{Results: The Capacity Trilogy}
\label{sec:results}

\begin{figure*}[t]
  \centering
  \includegraphics[width=0.98\textwidth]{fig_cliff}
  \caption{\textbf{Left:} at $d_{\mathrm{state}}{=}64$, held-out $h4$
  recovery transitions sharply as $K/\dstate$ increases; a logistic fit to
  six points (two archived anchors, $K{=}32/48$, plus four new intermediate
  loads, $K{=}34/38/42/46$) locates the midpoint at $x_0=0.5455$ (95\% CI
  $[0.5385,0.5513]$). \textbf{Right:} re-running the identical $K/\dstate$
  window at $d_{\mathrm{state}}{=}128$ ($K{=}68/76/84/92$) shows no
  transition anywhere in the window --- $h4{=}1.0$ at all 12 cells. The
  sigmoid fit on this flat data is degenerate by construction (bootstrap
  degenerate fraction $=1.0$) and is plotted as such, never as a spurious
  curve.}
  \label{fig:cliff}
\end{figure*}

\subsection{Act 1: the cliff, located}
\begin{itemize}
  \item At $d_{\mathrm{state}}=64$, six points spanning $K/\dstate \in
    [0.25, 0.75]$ (Fig.~\ref{fig:cliff}, left) trace a sharp, not gradual,
    transition: $h4 = 1.000$ ($K{=}16$) $\to 0.6669$ ($K{=}32$) $\to
    0.5676$ ($K{=}34$) $\to 0.3316$ ($K{=}38$) $\to 0.1177$ ($K{=}42) \to
    0.0434$ ($K{=}46$) $\to 0.0215$ ($K{=}48$).
  \item A logistic fit (0/4000 degenerate bootstrap fits) locates the
    transition's midpoint at $x_0 = 0.5455$ (95\% CI $[0.5385, 0.5513]$,
    width $0.0127$) with characteristic width $w=0.0597$ (CI $[0.0557,
    0.0642]$) --- i.e.\ the transition's effective span is $K/\dstate
    \approx [0.486, 0.605]$, or $K \approx 31$--$39$ raw at $d{=}64$. This
    is a $\sim\!20\times$ tightening of the prior 0.25-wide bracket
    (spanned only by the two flanking archived anchors) to a $\pm0.006$-wide
    located transition.
  \item \textbf{This is a WHERE result, not a WHY result.} The fit
    characterizes the transition's location and width; it makes no claim
    about the mechanism driving it (Acts 2--3 test the leading candidate
    directly).
\end{itemize}

\subsection{Act 2: the cliff, dissolved}
\begin{itemize}
  \item Re-running the identical $K/\dstate$ window at
    $d_{\mathrm{state}}=128$, using the equivalent loads $K=68,76,84,92$,
    shows $h4=1.0$ at all four loads, all 3 seeds each, $12/12$ cells, no
    exception (Fig.~\ref{fig:cliff}, right). This is the same $K/\dstate$
    window that collapsed $0.667 \to 0.022$ at $d=64$.
  \item Per pre-registered outcome semantics this is
    \textbf{CONFIRM-SHIFTED in its strong form}: not merely a shifted
    $x_0$, but the transition exits the measured window entirely. At
    $d=128$, capacity extends past $K/\dstate=0.71875$ --- at least 92 of
    107 entities recover perfectly through held-out 4-hop composition, in
    the exact band where $d=64$ had already collapsed to near zero.
  \item Instrument-saturation is ruled out: hop-21 recovery at $d=128$
    ranges $0.9900$--$1.0$ across the 12 cells (not uniformly $1.0$), and
    effective rank tracks each $K$ almost exactly (e.g.\ $K{=}84 \to
    83.6$) --- the state is using its full permitted rank, not saturating
    against $\dstate$ itself.
  \item \textbf{This rules out ``$K/\dstate$ ratio alone determines the
    cliff.''} The exact ratio window that cliffs at $d=64$ stays
    completely flat at $d=128$.
\end{itemize}

\begin{figure}[t]
  \centering
  \includegraphics[width=\linewidth]{fig_dose}
  \caption{Coherence-dose response at $K{=}68$, $d_{\mathrm{state}}{=}128$,
  rank-4 (concentrated) injection: $h4$ stays flat at $1.0$ at every dose
  tested (green), including $0.40$ --- exceeding $d_{\mathrm{state}}{=}64$'s
  own measured trained-coherence band (grey, $[0.373,0.385]$) --- while
  $d_{\mathrm{state}}{=}64$'s own cliff (orange) collapses across the
  identical coherence range at the matched $K/\dstate$ window.}
  \label{fig:dose}
\end{figure}

\subsection{Act 3: the leading confound, tested and exonerated}
\begin{itemize}
  \item The two waves above differ in more than raw scale: the $d=64$
    table (107 entities in 64 dimensions) is forced non-orthogonal by the
    Welch bound (measured coherence $0.2842$ at construction, drifting to a
    final-checkpoint range-of-$K$-means of $0.373$--$0.385$ under normal
    training); the $d=128$ table (107 entities in 128 dimensions) is
    exactly orthogonal at construction ($\max|\cos|=0.000000$). This is the
    single most natural remaining mechanistic candidate for the cliff's
    disappearance.
  \item We test it directly: holding $\dstate=128$, $n=107$, $K=68$ fixed
    at exactly the geometry Act 2 measured flat, we inject controlled,
    \emph{frozen} coherence into the anchor table under a concentrated
    (rank-4) structure and sweep the dose --- $0.130$, $0.284$, $0.40$ ---
    the last exceeding $d=64$'s own realized training band
    (Fig.~\ref{fig:dose}).
  \item \textbf{Result: $h4=1.0$ at every cell, every dose, $10/10$, no
    exception} --- flat across the entire dose range, including $0.40$.
    This is the strongest possible exoneration this design can produce for
    the rank-4 structure: directly-injected anchor-table coherence, held
    frozen for the full run and verified never to drift (bit-identical
    across all 10 recorded checkpoints per cell, zero slack), does not
    reproduce the $d=64$ cliff at matched $K/\dstate/n$ geometry.
  \item Frozen-constancy and instrument-saturation are both independently
    verified: \texttt{max\_abs\_cos} is exactly bit-identical across every
    checkpoint at every dose (the frozen-table mechanism worked as built),
    and hop-21 recovery ranges $0.9987$--$1.0$ across all 10 cells (not
    uniformly $1.0$) --- the flat $h4=1.0$ is a real ceiling for this
    quantity, not a broken or saturated readout.
\end{itemize}
